\chapter{РЕАЛИЗАЦИЯ КЛИЕНТСКОЙ ЧАСТИ, ТЕСТИРОВАНИЕ И РАЗВЁРТЫВАНИЕ ВЕБ-ПРИЛОЖЕНИЯ}

Третья глава посвящена реализации клиентской части веб-приложения «Палантир», описанию картографической визуализации пространственно-временных данных, а также проверке работоспособности основных компонентов и подготовке системы к эксплуатации в составе серверной среды. Если в предыдущей главе основное внимание уделялось проектированию архитектуры, базы данных, слоя доступа к базе данных и серверного API, то в рамках данной главы рассматриваются вопросы фронтенда, интеграционного тестирования, подготовки к развёртыванию и проверки приложения после размещения~\cite{sommerville2015software}.

Поскольку разрабатываемая система предназначена для хранения, обработки и визуализации данных об исторических военных конфликтах, клиентская часть должна обеспечивать наглядное представление пространственно-временных данных на карте и корректное взаимодействие с серверным API. Практическая значимость данного этапа заключается в подтверждении того, что разработанные компоненты способны функционировать как части единой веб-информационной системы, объединяющей интерфейс пользователя, серверную логику и базу данных.

\section{Реализация клиентской части веб-приложения}

Клиентская часть веб-приложения «Палантир» предназначена для взаимодействия пользователя с системой и представления данных, полученных от серверного API. Её основная роль заключается не в хранении или обработке предметной информации, а в обеспечении навигации, отображении сведений и визуализации пространственно-временных объектов на карте.

В составе клиентского уровня предусматриваются элементы интерфейса, необходимые для работы с историческими военными конфликтами: область карты, панель выбора отображаемых слоёв, элементы временной навигации, средства выбора войны или операции, а также информационная область для отображения сведений о выбранном объекте. Данные для этих элементов поступают от серверной части в форматах JSON и GeoJSON. JSON используется для обычных справочных сведений, а GeoJSON --- для объектов, имеющих координаты или геометрию.

Клиентский уровень формирует запросы к API на основании действий пользователя. При выборе конфликта, операции, даты или слоя отображения клиент отправляет HTTP-запрос с соответствующими параметрами. Полученный ответ используется для обновления интерфейса: списков, карточек, элементов карты и состава отображаемых слоёв. Такой подход позволяет отделить визуальное представление от серверной логики и базы данных.

Основные элементы клиентской части и их связь с серверным API представлены в таблице~\ref{tab:client-elements}.

\begin{table}[H]
\caption{Элементы клиентской части веб-приложения}
\label{tab:client-elements}
\centering
\small
\begin{tabular}{|p{3.5cm}|p{5.4cm}|p{4.7cm}|}
\hline
Элемент клиентской части & Назначение & Связь с серверным API \\
\hline
Карта & Отображение событий, позиций соединений и зон контроля & Получение пространственных объектов в GeoJSON \\
\hline
Временная шкала & Выбор даты или периода для анализа хода конфликта & Передача временных параметров в запросах \\
\hline
Панель слоёв & Управление отображением армий, событий и зон контроля & Запрос данных по выбранным типам объектов \\
\hline
Информационная карточка & Отображение сведений о выбранном объекте & Получение детальной информации по идентификатору \\
\hline
\end{tabular}
\end{table}

Для реализации интерфейса могут использоваться HTML, CSS и JavaScript. CSS-фреймворк Bootstrap, уже рассматриваемый при описании клиентской части, может применяться для унификации внешнего вида элементов управления и ускорения разработки интерфейсных компонентов. Для фронтенд-реализации интерактивной карты дополнительно применяется библиотека Leaflet, обеспечивающая работу с картографической подложкой, маркерами, полигонами, полилиниями и группами слоёв. При этом использование Bootstrap и Leaflet не влияет на архитектурное разделение системы: клиентская часть остаётся потребителем API и не выполняет прямой доступ к базе данных.

Особое значение для клиентской части имеет корректная интерпретация данных, полученных от сервера. Справочные JSON-ответы используются для заполнения списков и информационных блоков, а GeoJSON-объекты преобразуются в элементы карты. При изменении временного параметра или набора слоёв клиентская часть должна обновлять отображаемые объекты, чтобы карта отражала выбранный временной срез исторического процесса.

Таким образом, клиентская часть выполняет роль визуального и навигационного уровня веб-приложения. Она связывает действия пользователя с серверным API и обеспечивает представление данных в форме, пригодной для анализа исторических военных событий. Наиболее важной частью данного представления является картографическая визуализация пространственно-временных данных.

\section{Реализация картографической визуализации операции}

Одной из ключевых частей клиентской стороны веб-информационной системы «Палантир» является модуль картографической визуализации операции. Данный модуль предназначен для отображения пространственных объектов, связанных с выбранной боевой операцией: событий, войсковых соединений, зон контроля и линии фронта. Реализация выполнена средствами HTML, JavaScript и библиотеки Leaflet, которая используется для создания интерактивной карты в веб-интерфейсе.

Страница карты операции содержит две основные области: боковую панель управления и область отображения карты. В боковой панели размещаются элементы навигации, сведения об операции, переключатели картографических слоёв, элемент временной шкалы, а также список объектов, отображаемых на карте. Основная область страницы предназначена непосредственно для вывода интерактивной карты. В HTML-разметке для этого используется контейнер с идентификатором \texttt{map}, к которому в дальнейшем подключается логика библиотеки Leaflet.

Подключение функциональной части страницы выполняется через JavaScript-файл \texttt{map-page.js}. В нём сосредоточена основная логика работы карты: получение параметров страницы, инициализация картографической подложки, загрузка сведений об операции, создание слоёв, добавление объектов на карту, управление отображением слоёв и навигация по объектам.

При открытии страницы карта получает идентификатор операции из строки запроса. Например, страница может быть открыта в формате, приведённом в листинге~\ref{code:map-url-example}.

\begin{code}
\captionof{listing}{Пример адреса страницы карты операции}
\label{code:map-url-example}
\vspace{-0.75cm}
\begin{minted}{text}
/map.html?operationId=1
\end{minted}
\end{code}

Для получения параметра используется функция, основанная на объекте \texttt{URLSearchParams}. Её фрагмент показан в листинге~\ref{code:map-query-param}.

\begin{code}
\captionof{listing}{Получение параметра из строки запроса}
\label{code:map-query-param}
\vspace{-0.75cm}
\begin{minted}{javascript}
function getQueryParam(name) {
    return new URLSearchParams(window.location.search).get(name);
}
\end{minted}
\end{code}

Полученный идентификатор сохраняется во внутреннем состоянии страницы. Это позволяет использовать одну страницу карты для разных операций, так как конкретная операция определяется не отдельным HTML-файлом, а параметром \texttt{operationId}.

Для хранения текущего состояния страницы используется объект \texttt{pageState}, структура которого представлена в листинге~\ref{code:map-page-state}.

\begin{code}
\captionof{listing}{Объект состояния страницы карты}
\label{code:map-page-state}
\vspace{-0.75cm}
\begin{minted}{javascript}
const pageState = {
    operationId: getQueryParam("operationId"),
    operation: null,
    featureLayers: []
};
\end{minted}
\end{code}

Поле \texttt{operationId} хранит идентификатор выбранной операции, поле \texttt{operation} предназначено для хранения данных операции, полученных от серверного API, а массив \texttt{featureLayers} используется для сохранения всех добавленных на карту объектов. Наличие общего списка объектов позволяет выполнять автоматическое масштабирование карты по всем отображаемым элементам.

Инициализация карты выполняется с помощью библиотеки Leaflet. Для карты задаются координаты центра по умолчанию и начальный уровень масштабирования, как показано в листинге~\ref{code:leaflet-map-init}.

\begin{code}
\captionof{listing}{Инициализация карты Leaflet}
\label{code:leaflet-map-init}
\vspace{-0.75cm}
\begin{minted}{javascript}
const DEFAULT_CENTER = [51.75, 36.19];
const DEFAULT_ZOOM = 7;

const map = L.map("map", {
    zoomControl: true,
    attributionControl: true
}).setView(DEFAULT_CENTER, DEFAULT_ZOOM);
\end{minted}
\end{code}

После создания объекта карты подключается картографическая подложка. В качестве источника тайлов используется сервис Яндекс.Карт. API-ключ вынесен в отдельный конфигурационный файл, что позволяет отделить настройки доступа к внешнему картографическому сервису от основной логики страницы. Пример подключения тайлового слоя приведён в листинге~\ref{code:leaflet-yandex-tiles}.

\begin{code}
\captionof{listing}{Подключение картографической подложки}
\label{code:leaflet-yandex-tiles}
\vspace{-0.75cm}
\begin{minted}[mathescape=false]{javascript}
const yandexTilesApiKey = window.PALANTIR_MAP_CONFIG?.YANDEX_TILES_API_KEY;
const yandexTileUrl =
    "https://tiles.api-maps.yandex.ru/v1/tiles/" +
    "?x={x}&y={y}&z={z}&lang=ru_RU&l=map" +
    `&projection=web_mercator&apikey=${yandexTilesApiKey}`;

L.tileLayer(yandexTileUrl, {
    maxZoom: 19,
    tileSize: 256
}).addTo(map);
\end{minted}
\end{code}

Такой подход обеспечивает возможность централизованного изменения параметров подключения к картографическому сервису без изменения основного кода визуализации.

Для отображения разных типов данных на карте используются отдельные логические слои. Это позволяет независимо управлять видимостью событий, зон контроля, войсковых соединений и линии фронта. Структура слоёв приведена в листинге~\ref{code:leaflet-layer-groups}.

\begin{code}
\captionof{listing}{Создание логических слоёв карты}
\label{code:leaflet-layer-groups}
\vspace{-0.75cm}
\begin{minted}{javascript}
const mapLayers = {
    events: L.layerGroup().addTo(map),
    controlZones: L.layerGroup().addTo(map),
    armies: L.layerGroup().addTo(map),
    frontLines: L.layerGroup().addTo(map)
};
\end{minted}
\end{code}

Каждый слой представляет собой объект \texttt{L.layerGroup}. В него добавляются только объекты соответствующего типа. Например, события операции помещаются в слой \texttt{events}, зоны контроля --- в \texttt{controlZones}, войсковые соединения --- в \texttt{armies}, а линия фронта --- в \texttt{frontLines}. Такое разделение упрощает управление картой и делает интерфейс более гибким для пользователя.

Переключение слоёв реализуется через элементы управления в боковой панели. Для каждого переключателя назначается обработчик события \texttt{change}. Если переключатель включён, слой добавляется на карту, если выключен --- удаляется. Универсальная функция обработки переключателя представлена в листинге~\ref{code:leaflet-layer-toggle}.

\begin{code}
\captionof{listing}{Управление видимостью слоя карты}
\label{code:leaflet-layer-toggle}
\vspace{-0.75cm}
\begin{minted}{javascript}
function bindLayerToggle(checkbox, layer) {
    checkbox.addEventListener("change", function () {
        if (checkbox.checked) {
            layer.addTo(map);
        } else {
            map.removeLayer(layer);
        }
    });
}
\end{minted}
\end{code}

Данная функция является универсальной и используется для всех слоёв карты. Благодаря этому код не дублируется для каждого типа объектов, а управление слоями выполняется единообразно.

Загрузка сведений об операции выполняется асинхронно. После получения идентификатора операции клиентская часть отправляет запрос к серверному API. Полученные данные сохраняются в состоянии страницы и используются для заполнения информационного блока. Общая структура такой функции приведена в листинге~\ref{code:map-load-operation}.

\begin{code}
\captionof{listing}{Загрузка сведений об операции}
\label{code:map-load-operation}
\vspace{-0.75cm}
\begin{minted}{javascript}
async function loadOperation(operationId) {
    const operation = await fetchOperation(operationId);
    pageState.operation = operation;

    renderOperationInfo(operation);
    setMapCenterFromOperation(operation);
}
\end{minted}
\end{code}

Функция \texttt{renderOperationInfo} отвечает за вывод основной информации об операции: идентификатора, названия, дат начала и окончания, а также описания. При формировании HTML-содержимого используется функция экранирования данных, что снижает риск некорректной вставки пользовательского или внешнего текста в разметку страницы.

Если в данных операции присутствуют координаты центра, карта автоматически перемещается к соответствующей области. Для этого используется функция \texttt{setMapCenterFromOperation}, показанная в листинге~\ref{code:map-center-from-operation}.

\begin{code}
\captionof{listing}{Установка центра карты по данным операции}
\label{code:map-center-from-operation}
\vspace{-0.75cm}
\begin{minted}{javascript}
function setMapCenterFromOperation(operation) {
    const lat = operation.centerLat ?? operation.latitude ?? operation.lat;
    const lng = operation.centerLng ?? operation.longitude ?? operation.lng;

    if (isValidCoordinate(lat) && isValidCoordinate(lng)) {
        map.setView([Number(lat), Number(lng)], operation.zoom ?? DEFAULT_ZOOM);
    }
}
\end{minted}
\end{code}

В данном фрагменте предусмотрена проверка корректности координат. Это позволяет избежать ошибки при попытке переместить карту к неопределённым или некорректным значениям.

Отображение объектов на карте выполняется с использованием разных типов графических элементов Leaflet. События и войсковые соединения отображаются маркерами, линия фронта --- полилинией, зоны контроля --- полигонами. Для каждого объекта создаётся всплывающее окно с краткой информацией.

Пример добавления события на карту приведён в листинге~\ref{code:map-event-marker}.

\begin{code}
\captionof{listing}{Добавление события на карту}
\label{code:map-event-marker}
\vspace{-0.75cm}
\begin{minted}[mathescape=false]{javascript}
function addEventMarker(event) {
    const marker = L.marker([event.lat, event.lng], { icon: eventIcon })
        .bindPopup(`
            <strong>${escapeHtml(event.title)}</strong><br>
            <span>${escapeHtml(formatDate(event.date))}</span><br>
            <span>${escapeHtml(event.description)}</span>
        `);

    marker.addTo(mapLayers.events);
    pageState.featureLayers.push(marker);
}
\end{minted}
\end{code}

В данном случае маркер создаётся по координатам события. К нему привязывается всплывающее окно, содержащее название события, дату и описание. После создания маркер добавляется в слой событий и сохраняется в массиве \texttt{featureLayers}.

Войсковые соединения отображаются схожим образом, однако помещаются в отдельный слой. Пример такой функции приведён в листинге~\ref{code:map-army-marker}.

\begin{code}
\captionof{listing}{Добавление войскового соединения на карту}
\label{code:map-army-marker}
\vspace{-0.75cm}
\begin{minted}[mathescape=false]{javascript}
function addArmyMarker(army) {
    const marker = L.marker([army.lat, army.lng], { icon: armyIcon })
        .bindPopup(`
            <strong>${escapeHtml(army.title)}</strong><br>
            <span>${escapeHtml(army.description)}</span>
        `);

    marker.addTo(mapLayers.armies);
    pageState.featureLayers.push(marker);
}
\end{minted}
\end{code}

Такое разделение позволяет пользователю при необходимости скрыть соединения, не затрагивая отображение других объектов карты.

Линия фронта отображается с помощью полилинии. Для её построения используется массив координат, последовательно соединяемых на карте. Пример приведён в листинге~\ref{code:map-front-line}.

\begin{code}
\captionof{listing}{Добавление линии фронта на карту}
\label{code:map-front-line}
\vspace{-0.75cm}
\begin{minted}[mathescape=false]{javascript}
function addFrontLine(frontLine) {
    const line = L.polyline(frontLine.coordinates, {
        weight: 4,
        dashArray: "8 6"
    }).bindPopup(`<strong>${escapeHtml(frontLine.title)}</strong>`);

    line.addTo(mapLayers.frontLines);
    pageState.featureLayers.push(line);
}
\end{minted}
\end{code}

Использование полилинии подходит для отображения протяжённых объектов, таких как линия фронта, маршрут движения соединений или направление наступления.

Зоны контроля отображаются в виде полигонов. Такой способ подходит для представления территорий, находящихся под контролем определённой стороны конфликта. Пример добавления полигона приведён в листинге~\ref{code:map-control-zone-render}.

\begin{code}
\captionof{listing}{Добавление зоны контроля на карту}
\label{code:map-control-zone-render}
\vspace{-0.75cm}
\begin{minted}[mathescape=false]{javascript}
function addControlZone(zone) {
    const polygon = L.polygon(zone.coordinates)
        .bindPopup(`
            <strong>${escapeHtml(zone.title)}</strong><br>
            <span>${escapeHtml(zone.description)}</span>
        `);

    polygon.addTo(mapLayers.controlZones);
    pageState.featureLayers.push(polygon);
}
\end{minted}
\end{code}

Полигон создаётся на основе массива координат, описывающего границы территории. После добавления на карту объект также сохраняется в общем массиве \texttt{featureLayers}.

Помимо отображения объектов на карте, реализован список объектов в боковой панели. Для каждого события и войскового соединения создаётся отдельный элемент списка. При нажатии на элемент карта перемещается к координатам выбранного объекта, как показано в листинге~\ref{code:map-object-navigation}.

\begin{code}
\captionof{listing}{Переход к объекту из боковой панели}
\label{code:map-object-navigation}
\vspace{-0.75cm}
\begin{minted}{javascript}
button.addEventListener("click", function () {
    map.setView(item.coordinates, 10);
});
\end{minted}
\end{code}

Это решение упрощает навигацию по карте, так как пользователь может перейти к нужному объекту не только вручную, но и через список в интерфейсе.

Для отображения всех объектов на карте используется функция автоматического масштабирования. Она формирует группу из всех добавленных объектов и подбирает границы отображения таким образом, чтобы объекты были видны пользователю. Пример функции представлен в листинге~\ref{code:map-fit-objects}.

\begin{code}
\captionof{listing}{Автоматическое масштабирование карты по объектам}
\label{code:map-fit-objects}
\vspace{-0.75cm}
\begin{minted}{javascript}
function fitMapToObjects() {
    if (pageState.featureLayers.length === 0) {
        map.setView(DEFAULT_CENTER, DEFAULT_ZOOM);
        return;
    }

    const group = L.featureGroup(pageState.featureLayers);
    map.fitBounds(group.getBounds(), {
        padding: [40, 40],
        maxZoom: 10
    });
}
\end{minted}
\end{code}

Если объекты отсутствуют, карта возвращается к центру и масштабу по умолчанию. Если объекты присутствуют, Leaflet вычисляет их общие границы и изменяет масштаб карты.

Для корректного вывода данных также используются вспомогательные функции. Функция \texttt{formatDate} приводит дату к локализованному формату, функция \texttt{isValidCoordinate} проверяет корректность координат, а функция \texttt{escapeHtml} выполняет экранирование специальных символов. Фрагмент функции экранирования приведён в листинге~\ref{code:map-escape-html}.

\begin{code}
\captionof{listing}{Экранирование текстовых данных для HTML-разметки}
\label{code:map-escape-html}
\vspace{-0.75cm}
\begin{minted}{javascript}
function escapeHtml(value) {
    return String(value)
        .replaceAll("&", "&amp;")
        .replaceAll("<", "&lt;")
        .replaceAll(">", "&gt;")
        .replaceAll('"', "&quot;")
        .replaceAll("'", "&#039;");
}
\end{minted}
\end{code}

Экранирование необходимо для безопасного вывода текстовых данных внутри HTML-разметки, особенно при формировании всплывающих окон и информационных блоков.

Таким образом, клиентская реализация карты операции включает инициализацию интерактивной карты, подключение картографической подложки, получение данных операции через API, распределение объектов по слоям, отображение маркеров, линий и полигонов, а также средства навигации по объектам. Реализованная структура позволяет представить боевую операцию как совокупность взаимосвязанных пространственных объектов и обеспечивает удобное визуальное восприятие данных в рамках веб-информационной системы «Палантир».

\section{Unit-тестирование серверной части приложения}

Для проверки корректности работы серверной части системы «Палантир» были разработаны unit-тесты. Данный вид тестирования позволяет проверить отдельные компоненты приложения изолированно от остальных частей системы. В рамках серверной части это особенно важно, поскольку бизнес-логика приложения распределена между контроллерами, сервисами и репозиториями.

Unit-тестирование выполнялось для компонентов, отвечающих за получение данных о военных конфликтах и связанных с ними сущностях. Основное внимание было уделено проверке корректности работы сервисного слоя и контроллеров API. Такой подход позволяет убедиться, что приложение правильно обрабатывает как успешные сценарии выполнения запроса, так и ситуации, при которых запрашиваемые данные отсутствуют.

Для реализации тестов использовались библиотека \texttt{xUnit} и механизм mock-объектов. Mock-объекты применяются для замены реальных зависимостей тестируемого компонента. Например, при тестировании сервиса вместо обращения к настоящему репозиторию используется имитация репозитория, заранее настроенная на возврат определённых данных. Благодаря этому тестирование выполняется без подключения к базе данных, что делает проверку быстрее и позволяет сосредоточиться именно на логике работы тестируемого класса.

Структура unit-тестов построена по принципу \texttt{Arrange} --- \texttt{Act} --- \texttt{Assert}. На этапе \texttt{Arrange} выполняется подготовка тестовых данных и настройка mock-объектов. На этапе \texttt{Act} вызывается тестируемый метод. На этапе \texttt{Assert} проверяется, соответствует ли полученный результат ожидаемому поведению.

В качестве примеров были реализованы следующие unit-тесты:
\begin{enumerate}
\item Проверка успешного получения информации о конфликте по идентификатору.
\item Проверка обработки ситуации, когда конфликт с указанным идентификатором отсутствует.
\item Проверка работы контроллера при получении списка театров военных действий по идентификатору конфликта.
\end{enumerate}

Первый тест проверяет корректность работы метода получения конфликта по \texttt{Id}. В тесте создаётся объект \texttt{War}, содержащий идентификатор, название и краткое описание конфликта. Далее mock-объект репозитория настраивается таким образом, чтобы при обращении к методу \texttt{GetByIdAsync(1)} возвращался подготовленный объект. После вызова метода сервиса выполняется проверка, что результат не равен \texttt{null}, а значения полей \texttt{Id}, \texttt{Title} и \texttt{Summary} соответствуют ожидаемым данным. Соответствующий unit-тест представлен в листинге~\ref{code:war-service-getbyid-success-test}.

\begin{code}
\captionof{listing}{Unit-тест успешного получения конфликта по идентификатору}
\label{code:war-service-getbyid-success-test}
\vspace{-0.75cm}
\begin{minted}{csharp}
[Fact]
public async Task GetByIdAsync_WhenWarExists_ReturnsWarResponse()
{
    // Arrange
    var war = new War
    {
        Id = 1,
        Title = "Вторая мировая война",
        Summary = "Крупнейший военный конфликт XX века"
    };

    var warRepositoryMock = new Mock<IWarRepository>();

    warRepositoryMock
        .Setup(repo => repo.GetByIdAsync(1))
        .ReturnsAsync(war);

    var warService = new WarService(warRepositoryMock.Object);

    // Act
    var result = await warService.GetByIdAsync(1);

    // Assert
    Assert.NotNull(result);
    Assert.Equal(1, result.Id);
    Assert.Equal("Вторая мировая война", result.Title);
    Assert.Equal("Крупнейший военный конфликт XX века", result.Summary);

    warRepositoryMock.Verify(repo => repo.GetByIdAsync(1), Times.Once);
}
\end{minted}
\end{code}

Дополнительно в данном тесте используется проверка \texttt{Verify}, которая подтверждает, что метод репозитория был вызван ровно один раз. Это позволяет контролировать не только возвращаемый результат, но и корректность взаимодействия сервиса с зависимостью.

Второй тест проверяет поведение сервиса в ситуации, когда запрашиваемый конфликт отсутствует в хранилище данных. Для этого mock-объект репозитория настраивается на возврат значения \texttt{null} при обращении к методу \texttt{GetByIdAsync(999)}. Ожидаемым результатом является возникновение исключения \texttt{KeyNotFoundException}. Такая проверка необходима для подтверждения того, что приложение корректно обрабатывает некорректные или несуществующие идентификаторы. Unit-тест негативного сценария приведён в листинге~\ref{code:war-service-getbyid-not-found-test}.

\begin{code}
\captionof{listing}{Unit-тест обработки отсутствующего конфликта}
\label{code:war-service-getbyid-not-found-test}
\vspace{-0.75cm}
\begin{minted}{csharp}
[Fact]
public async Task GetByIdAsync_WhenWarDoesNotExist_ThrowsKeyNotFoundException()
{
    // Arrange
    var warRepositoryMock = new Mock<IWarRepository>();

    warRepositoryMock
        .Setup(repo => repo.GetByIdAsync(999))
        .ReturnsAsync((War?)null);

    var warService = new WarService(warRepositoryMock.Object);

    // Act & Assert
    await Assert.ThrowsAsync<KeyNotFoundException>(
        () => warService.GetByIdAsync(999)
    );

    warRepositoryMock.Verify(repo => repo.GetByIdAsync(999), Times.Once);
}
\end{minted}
\end{code}

Данный тест относится к негативным сценариям. Его наличие важно, поскольку серверная часть должна не только возвращать данные при корректных запросах, но и предсказуемо реагировать на ситуации, когда данные не найдены.

Третий тест направлен на проверку контроллера, отвечающего за получение списка театров военных действий по идентификатору конфликта. В данном случае mock-объект создаётся уже для сервиса \texttt{ITheaterService}, а не для репозитория. Это соответствует роли контроллера в архитектуре приложения: контроллер не должен напрямую работать с базой данных, а должен обращаться к сервисному слою.

В тесте заранее формируется список объектов \texttt{TheaterResponse}, после чего mock-объект сервиса настраивается на возврат данного списка при вызове метода \texttt{GetByWarIdAsync(1)}. Далее вызывается метод контроллера \texttt{GetByWarId(1)}, а результат проверяется на соответствие типу \texttt{OkObjectResult}. Фрагмент теста контроллера приведён в листинге~\ref{code:theaters-controller-getbywarid-test}.

\begin{code}
\captionof{listing}{Unit-тест контроллера получения театров военных действий}
\label{code:theaters-controller-getbywarid-test}
\vspace{-0.75cm}
\begin{minted}{csharp}
[Fact]
public async Task GetByWarId_WhenTheatersExist_ReturnsOkResult()
{
    // Arrange
    var theaterResponses = new List<TheaterResponse>
    {
        new TheaterResponse
        {
            Id = 1,
            WarId = 1,
            Title = "Восточный фронт",
            Description = "Основной театр боевых действий"
        },
        new TheaterResponse
        {
            Id = 2,
            WarId = 1,
            Title = "Западный фронт",
            Description = "Театр боевых действий в Западной Европе"
        }
    };

    var theaterServiceMock = new Mock<ITheaterService>();

    theaterServiceMock
        .Setup(service => service.GetByWarIdAsync(1))
        .ReturnsAsync(theaterResponses);

    var controller = new TheatersController(theaterServiceMock.Object);

    // Act
    var result = await controller.GetByWarId(1);

    // Assert
    var okResult = Assert.IsType<OkObjectResult>(result);
    var returnedTheaters = Assert.IsType<List<TheaterResponse>>(okResult.Value);

    Assert.Equal(2, returnedTheaters.Count);
    Assert.Equal("Восточный фронт", returnedTheaters[0].Title);

    theaterServiceMock.Verify(service => service.GetByWarIdAsync(1), Times.Once);
}
\end{minted}
\end{code}

Данный тест подтверждает, что контроллер корректно обрабатывает результат, полученный от сервисного слоя, и возвращает HTTP-ответ с кодом успешного выполнения. Также проверяется, что в теле ответа содержится список театров военных действий, соответствующий подготовленным тестовым данным.

Таким образом, разработанные unit-тесты позволяют проверить ключевые сценарии работы серверной части приложения: получение данных, обработку отсутствующих записей и формирование ответа контроллером API. Использование mock-объектов обеспечивает изоляцию тестируемых компонентов от базы данных и других внешних зависимостей. Это повышает надёжность проверки и позволяет быстрее выявлять ошибки в логике работы сервисов и контроллеров.

\section{Тестирование взаимодействия клиентской, серверной части и базы данных}

Тестирование взаимодействия клиентской, серверной части и базы данных направлено на оценку корректной работы приложения в условиях обмена данными между его основными компонентами. В данном случае проверяется не отдельный программный модуль, а полный маршрут обработки пользовательского действия: от формирования запроса в интерфейсе до получения результата из базы данных и его последующего отображения на странице или на карте.

Для веб-приложения «Палантир» интеграционная проверка имеет принципиальное значение, поскольку система объединяет несколько функционально различных подсистем. Клиентская часть отвечает за визуализацию данных, настройку отображаемых слоёв и работу с временными параметрами. Серверная часть реализует обработку HTTP-запросов, прикладную логику и доступ к данным. База данных PostgreSQL/PostGIS обеспечивает хранение как справочной, так и пространственной информации, связанной с войнами, театрами военных действий, операциями, сторонами конфликта, событиями, положением соединений и зонами контроля~\cite{postgresql_docs,postgis_manual}.

В рамках интеграционной проверки основное внимание уделяется следующим аспектам:
\begin{itemize}
\item Корректности формирования и отправки запросов из клиентской части.
\item Доступности конечных точек серверного API и правильности структуры возвращаемых ответов.
\item Корректности обращения серверной части к базе данных при выборке и обработке информации.
\item Согласованности передачи справочных данных в формате JSON и пространственных данных в формате GeoJSON.
\item Правильности отображения полученных результатов в элементах интерфейса и на карте.
\end{itemize}

Проверка строится на основе типовых пользовательских сценариев, соответствующих предметной области приложения. К числу таких сценариев относятся загрузка перечня конфликтов и операций, получение сведений о выбранной сущности, запрос пространственных объектов для заданной даты, обновление отображаемых данных при изменении временного параметра, а также отображение событий и зон контроля в картографическом интерфейсе. Такой подход позволяет оценить не только доступность отдельных функций, но и связность всей системы в процессе её практического использования.

Состав основных сценариев интеграционной проверки представлен в таблице~\ref{tab:integration-test-scenarios}.

\begin{table}[H]
\caption{Сценарии интеграционной проверки веб-приложения}
\label{tab:integration-test-scenarios}
\centering
\scriptsize
\begin{tabular}{|p{3.5cm}|p{3.1cm}|p{4.6cm}|p{2.3cm}|}
\hline
Проверяемый сценарий & Компоненты & Ожидаемый результат & Комментарий \\
\hline
Получение списка сущностей & API, EF Core, БД & Возвращается JSON-ответ с набором записей & Справочные данные \\
\hline
Получение записи по идентификатору & API, репозиторий, БД & Возвращается объект или корректный ответ при отсутствии данных & Детальная карточка \\
\hline
Получение пространственных данных & API, PostGIS, клиентская часть & Возвращается структура, пригодная для отображения на карте & GeoJSON-данные \\
\hline
Передача временного параметра & Клиент, API, БД & Набор объектов соответствует выбранной дате или периоду & Временная навигация \\
\hline
Проверка отображения данных & Клиентская часть, API & Полученные сведения используются для обновления интерфейса & Сквозной сценарий \\
\hline
\end{tabular}
\end{table}

На уровне взаимодействия клиентской и серверной частей проверяется корректность передачи параметров запросов, в том числе идентификаторов сущностей, фильтров и временных характеристик. Существенное значение имеет проверка того, что сервер принимает корректно сформированные запросы, выполняет необходимую обработку и возвращает ответы в формате, пригодном для использования клиентским интерфейсом. Для этого анализируется структура ответов API, наличие необходимых полей в полезной нагрузке, соответствие структуры данных ожидаемой модели и корректность обработки ситуаций, связанных с отсутствием записей или передачей неполных параметров~\cite{aspnet_docs}.

На уровне серверной части и базы данных проверяется согласованность прикладной модели и структуры хранения. Обращение к данным должно обеспечивать получение непротиворечивой информации о временных и пространственных объектах. Для приложения данного типа особенно важно, чтобы выборка геометрических данных сопровождалась корректной передачей временной привязки, идентификаторов операций и принадлежности объектов к сторонам конфликта.

Таким образом, тестирование взаимодействия клиентской, серверной части и базы данных позволяет оценить согласованность основных компонентов веб-приложения «Палантир». Рассматриваемые сценарии показывают, каким образом архитектура обеспечивает передачу данных между уровнями системы и поддерживает решение задач хранения и визуализации сведений об исторических военных конфликтах.

\section{Развёртывание веб-приложения на сервере}

После завершения основных этапов реализации приложения следующим шагом является его подготовка к размещению в серверной среде. Развёртывание необходимо для перевода системы из локального режима разработки в режим, при котором доступ к её функциональности может осуществляться через сетевое взаимодействие между клиентской и серверной частями. На данном этапе особое значение имеет согласованность конфигурации всех компонентов приложения и обеспечение их совместной работы.

С учётом выбранной архитектуры подготовка к развёртыванию включает публикацию серверной части, подготовку параметров подключения к базе данных, размещение клиентского интерфейса и настройку взаимодействия между этими компонентами. Серверная часть, реализованная на платформе ASP.NET Core, должна быть размещена в среде, обеспечивающей её запуск и обработку HTTP-запросов. База данных PostgreSQL/PostGIS должна быть доступна серверному приложению по корректно настроенным параметрам подключения. Клиентская часть, в свою очередь, должна использовать актуальный адрес серверного API и обеспечивать корректную загрузку интерфейса и картографических данных~\cite{aspnet_configuration,aspnet_docs}.

В общем виде процесс подготовки к развёртыванию включает следующие этапы:
\begin{itemize}
\item Подготовку конфигурационных параметров приложения, включая параметры доступа к базе данных и адреса взаимодействия между компонентами.
\item Публикацию серверной части в выбранную среду выполнения.
\item Подготовку базы данных и обеспечение доступности необходимой структуры хранения для серверного приложения.
\item Размещение клиентской части и настройку её взаимодействия с опубликованным API.
\item Контрольный запуск приложения и первичную проверку доступности его основных компонентов.
\end{itemize}

При развёртывании существенное значение имеет согласованность конфигурации среды. Даже при корректной локальной реализации приложение не сможет функционировать в серверном режиме при некорректных параметрах подключения к базе данных, несоответствии адресов API или отсутствии доступа к необходимым ресурсам. По этой причине на этапе развёртывания необходимо обеспечить единство настроек, относящихся к соединению с СУБД, маршрутизации запросов и обработке клиентских обращений.

Отдельное внимание уделяется вопросу переноса данных и структуры базы. Для рассматриваемого приложения это особенно важно, поскольку в системе используются не только текстовые и справочные сведения, но и пространственные объекты. Следовательно, серверная среда должна обеспечивать доступность расширения PostGIS и возможность корректной обработки геометрических данных. При наличии предварительно подготовленной структуры БД необходимо также удостовериться, что используемая схема соответствует прикладной модели, реализованной в серверной части приложения~\cite{postgis_manual,efcore_scaffolding}.

Размещение клиентской части завершает процесс подготовки веб-приложения к работе на сервере. После этого приложение получает возможность функционировать как единая система, в которой пользовательский интерфейс инициирует запросы, серверная часть выполняет необходимую обработку, а база данных обеспечивает выдачу требуемых сведений для последующего отображения. Тем самым развёртывание выступает не формальной операцией публикации файлов, а практическим этапом интеграции всех реализованных компонентов.

\section{Проверка работоспособности приложения после развёртывания на сервере}

Для этапа проверки после размещения приложения в серверной среде необходимо определить набор контрольных действий, позволяющих оценить его работоспособность. Цель данного этапа заключается в подтверждении того, что система сохраняет корректное поведение не только в условиях локальной разработки, но и после переноса в среду эксплуатации. Такая проверка позволяет выявить возможные проблемы конфигурации, сетевого взаимодействия, подключения к базе данных и отображения интерфейса, которые не всегда проявляются на предыдущих этапах.

Постразвёрточная проверка должна охватывать как доступность отдельных компонентов, так и выполнение сквозных пользовательских сценариев. В первую очередь оценивается доступность серверного API и его готовность к обработке запросов. Далее подтверждается корректность подключения к базе данных и возможность получения из неё как справочной, так и пространственной информации. После этого проверяется клиентская часть: корректность загрузки страниц, работа элементов интерфейса, отправка запросов к API и отображение полученных результатов.

В ходе такой проверки целесообразно рассматривать следующие сценарии:
\begin{itemize}
\item Открытие клиентского интерфейса и загрузку основных элементов страницы.
\item Получение справочных данных о конфликтах, театрах военных действий и операциях.
\item Отображение пространственных объектов на карте после обращения к серверному API.
\item Обновление данных при изменении временных параметров и фильтров отображения.
\item Проверку базовых пользовательских сценариев навигации и просмотра сведений о выбранных сущностях.
\end{itemize}

Для системы «Палантир» особое значение имеет подтверждение того, что после развёртывания сохраняется корректность визуализации исторических военных данных на карте. Это означает, что пользователь должен иметь возможность получать сведения об операциях, событиях, соединениях и зонах контроля в форме, согласованной с временным контекстом и пространственным положением объектов. Если после публикации приложения такие сценарии выполняются без нарушения структуры данных и без рассогласования между интерфейсом, серверной частью и базой данных, то это является основанием для вывода о корректной работоспособности системы в серверной среде.

Итоговая проверка после размещения приложения позволяет оценить готовность основных компонентов к совместной работе и определить направления дальнейшего уточнения конфигурации. Рассмотренные действия по тестированию, подготовке к развёртыванию и постразвёрточному контролю показывают, что приложение «Палантир» рассматривается как целостное программное средство для хранения, обработки и визуализации данных об исторических военных конфликтах.
